% ---
% titre: Problème de proportionnalité - le remplissage de bassin
% theme: FA
% chapitre: Proportionnalité
% sous_chapitre: Résoudre des problèmes
% niveau: [10e]
% type: EVAL
% tags: [proportionnalite, p-question-TS]
% difficulte: 3
% corrige: true
% ---

\question[4]
Un robinet qui débite 18 litres à la minute met 28 heures pour remplir un bassin. Quel temps mettrait-il si son débit était de 42 litres à la minute ?

\vspace{10pt}{\footnotesize\raggedleft Espace pour ta démarche et tes calculs\par}
\begin{solutionorgrid}[300pt]

\vspace{5pt}
Conversion d'unité de temps: 28 heures = $28 \cdot 60 = 1680$ minutes \hfill\textbf{(0,5pt)}

\vspace{5pt}
\def\arraystretch{1.5}
\begin{tabularx}{0.5\textwidth}{l|C|C|}
Temps [min]&1&1680\\
\hline
Volume [l]&18&\\
\end{tabularx}
\def\arraystretch{1}

\hfill\textbf{(0,5pt)} pour la construction du tableau

\vspace{15pt}$18 \cdot 1680 : 1 = 30\,240$\,l \hfill calcul \textbf{(0,5pt)} et réponse \textbf{(0,5pt)}

\vspace{20pt}
\def\arraystretch{1.5}
\begin{tabularx}{0.5\textwidth}{l|C|C|}
Temps [min]&1&\\
\hline
Volume [l]&42&30\,240\\
\end{tabularx}
\def\arraystretch{1}

\hfill\textbf{(0,5pt)} pour la construction du tableau

\vspace{15pt}$30\,240 \cdot 1 : 42 = 720$\,min \hfill calcul \textbf{(0,5pt)} et réponse \textbf{(0,5pt)}

\vspace{5pt}
Conversion d'unité de temps: 720 minutes = $720 \cdot 60 = 12$ heures
\end{solutionorgrid}

\begin{tcolorbox}[enhanced jigsaw, interior hidden, colframe=box, parbox=false]%
Réponse: 
\sol{\vspace{20pt}}{Le bassin mettrai 12 heures à se remplir. \textbf{(0,5pt)}}
\end{tcolorbox}



